% !TEX ROOT = ../ersti.tex
\section{Studium im Ausland}
% Man hört immer wieder von den tollen Erfahrungen und Möglichkeiten, die ein Studium im Ausland bietet. Die Leute schwärmen davon, wie schön es an der Uni so-und-so war. Aber wie kommt man überhaupt dahin?

% Auf der Homepage der jeweiligen Fakultät findet ihr eine Liste von Unis, mit denen Austauschprogramme laufen, zum Beispiel über ERASMUS oder bilaterale Abkommen, die also üblicherwiese anerkannt werden, einige generelle Infos und ein paar Erfahrungsberichte, die ihr bei Interesse anschauen solltet. Die Gleichwertigkeit von Studienleistungen muss allerdings nicht von euch überprüft werden, die Beweispflicht liegt bei der Fakultät. Lasst euch also nicht beirren und hört im Zweifel die Veranstaltung im Ausland einfach, die sollen erst mal zeigen, dass das nicht das Gleiche war.

% Allerdings sollte ein Auslandssemester im Moment nicht euer erstes Problem sein -- während der ersten beiden Semester ist es ganz einfach fachlich nicht möglich. Bis zur Einführung der Bachelor- und Masterstudiengänge konnte man -- jedenfalls in der Mathe -- an diesen Programmen nur nach abgeschlossenem Grundstudium teilnehmen, inzwischen wird der Zeitraum letztes Bachelor Jahr\,/\,erstes Master Jahr empfohlen. Das bedeutet, dass ihr nach dem zweiten Semester damit beginnen solltet, euch umzuschauen, insbesondere in Bezug auf Bewerbungsvoraussetzungen und -fristen.

% Auslands-BAföG müsst ihr übrigens separat beantragen, das zu\-stän\-dige Amt dafür hängt vom jeweiligen Ausland ab\footnote{\url{https://www.auslandsbafoeg.de/}}. In Heidelberg kann man sich auch direkt an das Dezernat Internationale Beziehungen, früher Akademisches Auslandsamt, (AAA)\footnote{\url{https://www.uni-heidelberg.de/einrichtungen/verwaltung/internationales/}} wenden, das ihr in der Seminarstraße 2 findet. Ihr müsst dann ins Auslandsstudium Info-Zimmer 139\footnote{Öffnungszeiten: \auslandsinfooeff}.

% Sehr wichtig ist noch zu wissen, dass man Urlaubssemester beantragen kann, damit die Fachsemesterzahl -- vor allem wichtig für alle, die Inlands-BAföG erhalten -- nicht weiter läuft. Dafür läuft die Hochschulsemesterzahl weiter. Man kann sich dann die im Ausland erworbenen \gls{LP} -- nach Absprache mit dem Prüfungssekretariat -- anrechnen lassen.

% Für das Urlaubssemester ist die Studierendenadministration in der Seminarstraße 2 zuständig. Mehr Infos und das Formular zum Urlaubssemester findet ihr online\footnote{\url{https://www.uni-heidelberg.de/studium/imstudium/formalia/beurlaubung.html}}.
 
Man hört immer wieder von den tollen Erfahrungen und Möglichkeiten, die ein Studium im Ausland bietet. Die Leute schwärmen davon, wie schön es an der Uni so-und-so war. Aber wie kommt man überhaupt dahin? 

Auf der Homepage der jeweiligen Fakultät und der Universität Heidelberg findet ihr eine Liste von Unis, mit denen die Universität Heidelberg  Austauschprogramme hat, und die relevanten Informationen zu Erasmus und anderen Austauschprogrammen. Man profitiert in Heidelberg von einer Vielzahl an Austauschprogramme mit Universitäten aus der ganzen Welt, bei vielen dieser gibt es außerdem einen Studiengebührenerlass für Studenten aus Heidelberg, was für viele ausländische Universitäten mit hohen Semestergebühren ein großer Vorteil ist. 

Allerdings sollte ein Auslandssemester im Moment nicht euer erstes Problem sein -- während der ersten Semestern ist es ganz einfach fachlich nicht möglich oder sinnvoll, ein Auslandssemester zu machen. Nach abgeschlossenem Grundstudium, also den letzten beiden Bachelor Semester ist ein Auslandssemester zu empfohlen. Das bedeutet, dass ihr nach dem ersten Semester damit beginnen solltet, euch umzuschauen, insbesondere in Bezug auf Bewerbungsvoraussetzungen und -fristen. Die Bewerbungsverfahren sind je nach Austauschprogramm sehr früh, teilweise über ein Jahr vor! Auf dieser Webseite findet ihr die jeweiligen Austauschprogramme mit den dazugehörigen fristen \footnote{\url{https://www.uni-heidelberg.de/de/studium/studium-international/studium-im-ausland/austauschprogramme-weltweit}}.
In Heidelberg kann man sich auch direkt an das Infocenter das Dezernat internationale Beziehungen wenden \footnote{\url{https://www.uni-heidelberg.de/de/studium/studium-international/studium-im-ausland/infocenter-studium-im-ausland}}. Diese bieten auch Sprechstunden an, bei denen man sich gut informieren kann.

Auslands-BAföG müsst ihr übrigens separat beantragen, das zu\-stän\-dige Amt dafür hängt vom jeweiligen Ausland ab\footnote{\url{https://www.auslandsbafoeg.de/}}. 

Sehr wichtig ist noch zu wissen, dass man Urlaubssemester beantragen kann, damit die Fachsemesterzahl -- vor allem wichtig für alle, die Inlands-BAföG erhalten -- nicht weiter läuft. Dafür läuft die Hochschulsemesterzahl weiter. Für das Urlaubssemester ist die Studierendenadministration in der Seminarstraße 2 zuständig. Mehr Infos und das Formular zum Urlaubssemester findet ihr online\footnote{\url{https://www.uni-heidelberg.de/studium/imstudium/formalia/beurlaubung.html}}. 

Im Allgemeinen könnt ihr eure Studienleistungen aus dem Ausland anrechnen lassen, es ist sinnvoll vor der Wahl der Moodle sich mit dem Erasmus Fachkoordinator eures jeweiligen Fachs auszutauschen, das ist die jeweilige Ansprechperson, auch wenn nicht über Erasmus ins Ausland geht. Für fachspezifische fragen sind diese auch eine gute Anlaufstelle.
Es ist auch sinnvoll, wen ihr ein Auslandssemester vorhabt dies bei eurer Studiumsplan zu beachten.
